\subsection{渐近展开的和与积}

若 \(f_{1}\) 、 \(f_{2}\) 相对于一个比较尺度 E 分别容许精确到 \(g_{\alpha}\) 与 \(g_{\beta}\) 的渐近展开，则通过把两个展开限制到这个精确度，可导出精确到 \(g_{\min(\alpha,\beta)}\) 的展开；我们已在 V, p.~222 中看到如何得到 \(f_{1} + f_{2}\) 的一个精确到 \(g_{\min(\alpha,\beta)}\) 的渐近展开。

设 \(V_{1}, V_{2}\) 与 V 是域 K 上的三个赋范空间，并设 \((\mathbf{x}, \mathbf{y}) \mapsto [\mathbf{x}, \mathbf{y}]\) 是从 \(V_{1} \times V_{2}\) 到 V 中的一个连续双线性映射；对本节的其余部分，我们进一步假设尺度 E 使得 E 中任意两个函数的乘积仍在 E 中（这对所有作为例子给出的比较尺度（在 V, p.~220 中）都成立）。

现在设 \(f_{1}\) 、 \(f_{2}\) 分别是 \(\mathcal{H}(\mathfrak{F}, V_{1})\) 与 \(\mathcal{H}(\mathfrak{F}, V_{2})\) 中的两个函数，分别具有渐近展开 \(f_{1} = \sum_{\lambda \leqslant \alpha} a_{\lambda} g_{\lambda} + r_{\alpha}\) 与 \(f_{2} = \sum_{\mu \leqslant \beta} b_{\mu} g_{\mu} + r_{\beta}\) （精确度分别为 \(g_{\alpha}\) 与 \(g_{\beta}\) ，相对于尺度 E）。进一步假设 \(a_{\lambda}\) 与 \(b_{\mu}\) 都不全为零，并设 \(a_{\gamma} g_{\gamma}\) 与 \(b_{\delta} g_{\delta}\) 是 \(f_{1}\) 与 \(f_{2}\) 的主部。根据假设，可以写 \(g_{\gamma} g_{\beta} = g_{\rho}\) 与 \(g_{\delta} g_{\alpha} = g_{\sigma}\) ；让我们证明：和 \(\sum[a_{\lambda}.b_{\mu}]g_{\lambda}g_{\mu}\) （对所有对 \((\lambda,\mu)\) （满足 \(g_{\min(\rho,\sigma)}\ll g_{\lambda}g_{\mu}\) ）取的）是 \([f_{1}.f_{2}]\) 的一个精确到 \(g_{\min(\rho,\sigma)}\) 的渐近展开。现在 \([f_{1}.f_{2}]\) 与这个和的差是有限多个项之和，其中每一项或者形如 \([a_{\lambda}.b_{\mu}]g_{\lambda}g_{\mu}\) （满足 \(g_{\lambda}g_{\mu}\ll g_{\min(\rho,\sigma)}\) ），或者形如 \([a_{\lambda}.r_{\beta}]g_{\lambda}\) （其中 \(\lambda\geqslant\gamma\) ），或者形如 \([r_{\alpha}.b_{\mu}]g_{\mu}\) （其中 \(\mu\geqslant\delta\) ）；但由于 {[}x.y{]} 连续，由 (V, p.~213, prop. 3 and V, p.~215, prop. 6) 有 \([a_{\lambda}.r_{\beta}]g_{\lambda}\preccurlyeq r_{\beta}g_{\lambda}\ll g_{\beta}g_{\gamma}=g_{\rho}\) （对 \(\lambda\geqslant\gamma\) ），类似地 \([r_{\alpha}.b_{\mu}]g_{\mu}\preccurlyeq r_{\alpha}g_{\mu}\ll g_{\alpha}g_{\delta}=g_{\sigma}\) （对 \(\mu\geqslant\delta\) ），由此得命题 (V, p.~215, prop. 5)。

若所有 \(a_{\lambda}\) 为零，则 \([f_{1}.f_{2}] \ll g_{\alpha}g_{\delta}\) ；换句话说，有 \([f_{1}.f_{2}]\) 的一个零项的、精确到 \(g_{\alpha}g_{\delta}\) 的渐近展开；类似地，若所有 \(a_{\lambda}\) 与 \(b_{\mu}\) 为零，则有 \([f_{1}.f_{2}]\) 的一个零项的、精确到 \(g_{\alpha}g_{\beta}\) 的渐近展开。

我们将把上述结果主要应用于 V 是 K 上的赋范代数且双线性函数 \([x.y]\) 是这个代数中的乘积 xy 的情形；最重要的情形是 V 等于 R 或 C 的情形。

特别地，若 \(f_{i}\) （ \(1 \leqslant i \leqslant n\) ）是 \(\mathcal{H}(\mathfrak{F}, \mathrm{K})\) 中的 n 个函数，其中每个都相对于 E 容许一个渐近展开，则对每个多项式 \(\sum_{(v_{i})} a_{v_{1}v_{2}\ldots v_{n}} f_{1}^{v_{1}} \ldots f_{n}^{v_{n}}\) （ \(f_{i}\) 的，系数在赋范空间 V 中）都可以得到一个相对于 E 的渐近展开；此外，前面的规则使我们能确定所得到的展开的精确度，只要我们知道函数 \(f_{i}\) 的展开的精确度。

